% --- CORRECCION: Forzar interpretacion raw de la entrada ---
\UseRawInputEncoding
\documentclass[aspectratio=169, 10pt, xcolor=table]{beamer}

% --- Configuracion del Tema ---
\usetheme{Madrid}
\usecolortheme{dolphin}

% --- Paquetes necesarios ---
\usepackage[utf8]{inputenc}
\usepackage[T1]{fontenc}
\usepackage[spanish,es-noquoting]{babel}
\usepackage{amsmath, amssymb}
\usepackage{mathtools}
\usepackage{booktabs}
\usepackage{graphicx}
\usepackage{listings}
\usepackage{tikz}
\usetikzlibrary{arrows.meta, positioning, shapes.geometric, babel}
\usepackage{etoolbox}
\usepackage{upquote}

% --- Estilo para codigo SQL ---
\definecolor{codegreen}{rgb}{0,0.6,0}
\definecolor{codegray}{rgb}{0.5,0.5,0.5}
\definecolor{codepurple}{rgb}{0.58,0,0.82}
\definecolor{backcolour}{rgb}{0.95,0.95,0.92}
\lstset{
    language=SQL,
    backgroundcolor=\color{backcolour},
    commentstyle=\color{codegreen},
    keywordstyle=\color{blue},
    numberstyle=\tiny\color{codegray},
    stringstyle=\color{codepurple},
    basicstyle=\ttfamily\scriptsize,
    breaklines=true,
    frame=single,
    showstringspaces=false,
    literate={á}{{\'a}}1 {é}{{\'e}}1 {í}{{\'i}}1 {ó}{{\'o}}1 {ú}{{\'u}}1 {ñ}{{\~n}}1 {ü}{{\"u}}1
}

% --- Pie de pagina personalizado ---
\makeatletter
\setbeamertemplate{footline}{
  \leavevmode%
  \hbox{%
  \begin{beamercolorbox}[wd=.333333\paperwidth,ht=2.25ex,dp=1ex,center]{author in head/foot}%
    \usebeamerfont{author in head/foot}\insertshortauthor
  \end{beamercolorbox}%
  \begin{beamercolorbox}[wd=.333333\paperwidth,ht=2.25ex,dp=1ex,center]{title in head/foot}%
    \usebeamerfont{title in head/foot}Sesión 6
  \end{beamercolorbox}%
  \begin{beamercolorbox}[wd=.333333\paperwidth,ht=2.25ex,dp=1ex,right]{date in head/foot}%
    \usebeamerfont{date in head/foot}BBDD \hfill \insertframenumber/\inserttotalframenumber \hspace*{0.3cm}
  \end{beamercolorbox}}%
  \vskip0pt%
}
\makeatother

% --- Datos de la portada ---
\title[SESIÓN 6 BBDD]{\textbf{Sesión 6}}
\subtitle{Análisis Avanzado y Transformación de Datos con SQL}
\author[Prof. C. Sánchez]{Carlos César Sánchez Coronel}
\institute{\textbf{BASES DE DATOS}}
\date{}

\setbeamertemplate{navigation symbols}{}

\AtBeginSection[]{
  \begin{frame}
    \frametitle{Contenido de la sección}
    \tableofcontents[currentsection]
  \end{frame}
}

\begin{document}

\begin{frame}
    \titlepage
\end{frame}

\begin{frame}{Contenido general}
    \tableofcontents
\end{frame}

% ============================================================================
% SECCIÓN 1: AGREGACIONES Y AGRUPAMIENTOS
% ============================================================================
\section{Agregaciones y Agrupamientos}

\begin{frame}{Funciones de Agregación Básicas}
    \begin{block}{Funciones estándar}
        \begin{itemize}
            \item \texttt{COUNT(*)}: número de filas
            \item \texttt{COUNT(expr)}: valores no nulos
            \item \texttt{SUM(expr)}: suma
            \item \texttt{AVG(expr)}: promedio
            \item \texttt{MIN(expr)}: mínimo
            \item \texttt{MAX(expr)}: máximo
        \end{itemize}
    \end{block}
\end{frame}

\begin{frame}[fragile]{GROUP BY y HAVING}
    \begin{block}{Ejemplo: salario promedio por departamento}
        \begin{lstlisting}
SELECT departamento_id,
       COUNT(*) AS num_empleados,
       AVG(salario) AS salario_promedio
FROM empleados
GROUP BY departamento_id;
        \end{lstlisting}
    \end{block}
    \begin{block}{Filtrar grupos con HAVING}
        \begin{lstlisting}
SELECT departamento_id, AVG(salario) AS promedio
FROM empleados
GROUP BY departamento_id
HAVING AVG(salario) > 4000;
        \end{lstlisting}
    \end{block}
\end{frame}

\begin{frame}[fragile]{Agregaciones Avanzadas: ROLLUP, CUBE, GROUPING SETS}
    \begin{block}{ROLLUP – subtotales jerárquicos}
        \begin{lstlisting}
SELECT departamento_id, ciudad, SUM(salario)
FROM empleados
GROUP BY ROLLUP (departamento_id, ciudad);
        \end{lstlisting}
    \end{block}
    \begin{block}{CUBE – todas las combinaciones}
        \begin{lstlisting}
SELECT departamento_id, ciudad, SUM(salario)
FROM empleados
GROUP BY CUBE (departamento_id, ciudad);
        \end{lstlisting}
    \end{block}
    \begin{block}{GROUPING SETS – agrupaciones específicas}
        \begin{lstlisting}
SELECT departamento_id, ciudad, SUM(salario)
FROM empleados
GROUP BY GROUPING SETS ((departamento_id), (ciudad), ());
        \end{lstlisting}
    \end{block}
\end{frame}

\begin{frame}[fragile]{Ejercicio Resuelto: Resumen de Ventas}
    \textbf{Enunciado:} Tabla \texttt{ventas(producto, categoría, monto, fecha)}. Obtener:
    \begin{enumerate}
        \item Monto total por categoría.
        \item Monto total por categoría y producto.
        \item Categorías con total > 10000.
        \item Totales con ROLLUP.
    \end{enumerate}
    \begin{block}{Solución}
        \begin{lstlisting}
-- 1. Por categoría
SELECT categoria, SUM(monto) AS total
FROM ventas GROUP BY categoria;

-- 2. Por categoría y producto
SELECT categoria, producto, SUM(monto)
FROM ventas GROUP BY categoria, producto;

-- 3. Categorías con total > 10000
SELECT categoria, SUM(monto) AS total
FROM ventas
GROUP BY categoria
HAVING SUM(monto) > 10000;

-- 4. Con ROLLUP
SELECT categoria, producto, SUM(monto)
FROM ventas
GROUP BY ROLLUP (categoria, producto);
        \end{lstlisting}
    \end{block}
\end{frame}

% ============================================================================
% SECCIÓN 2: FUNCIONES DE VENTANA
% ============================================================================
\section{Funciones de Ventana}

\begin{frame}{Concepto y Sintaxis}
    \begin{block}{Definición}
        \begin{itemize}
            \item Calculan sobre un conjunto de filas relacionadas con la fila actual.
            \item No colapsan filas (mismo número de filas de salida).
            \item Se usan con \texttt{OVER()}.
        \end{itemize}
    \end{block}
    \begin{exampleblock}{Sintaxis básica}
        \texttt{funcion() OVER (PARTITION BY col1 ORDER BY col2)}
    \end{exampleblock}
\end{frame}

\begin{frame}[fragile]{Funciones de Ranking}
    \begin{block}{\texttt{ROW\_NUMBER()}, \texttt{RANK()}, \texttt{DENSE\_RANK()}}
        \begin{lstlisting}
SELECT nombre, salario,
       ROW_NUMBER() OVER (ORDER BY salario DESC) AS row_num,
       RANK()       OVER (ORDER BY salario DESC) AS rank,
       DENSE_RANK() OVER (ORDER BY salario DESC) AS dense_rank
FROM empleados;
        \end{lstlisting}
    \end{block}
    \begin{itemize}
        \item \texttt{ROW\_NUMBER}: secuencial único
        \item \texttt{RANK}: mismo rango para empates, deja huecos
        \item \texttt{DENSE\_RANK}: mismo rango, sin huecos
    \end{itemize}
\end{frame}

\begin{frame}[fragile]{Funciones de Valor: LAG y LEAD}
    \begin{block}{Acceder a filas anteriores/posteriores}
        \begin{lstlisting}
SELECT fecha, monto,
       LAG(monto, 1) OVER (ORDER BY fecha) AS monto_anterior,
       monto - LAG(monto, 1) OVER (ORDER BY fecha) AS diferencia
FROM ventas_diarias;
        \end{lstlisting}
    \end{block}
    \begin{block}{\texttt{FIRST\_VALUE}, \texttt{LAST\_VALUE}}
        \begin{lstlisting}
SELECT producto, fecha, cantidad,
       FIRST_VALUE(cantidad) OVER (PARTITION BY producto ORDER BY fecha) AS primera_venta
FROM ventas;
        \end{lstlisting}
    \end{block}
\end{frame}

\begin{frame}[fragile]{Agregaciones como Ventana y Marcos}
    \begin{block}{Totales acumulados y medias móviles}
        \begin{lstlisting}
SELECT fecha, monto,
       SUM(monto) OVER (ORDER BY fecha) AS acumulado,
       AVG(monto) OVER (ORDER BY fecha ROWS BETWEEN 2 PRECEDING AND CURRENT ROW) AS media_movil_3
FROM ventas_diarias;
        \end{lstlisting}
    \end{block}
    \begin{itemize}
        \item \texttt{ROWS BETWEEN ... AND ...} define el marco.
        \item \texttt{RANGE} trabaja con valores (ej. fechas).
    \end{itemize}
\end{frame}

\begin{frame}[fragile]{Ejemplo: Top N por Departamento}
    \begin{block}{Tres empleados mejor pagados por departamento}
        \begin{lstlisting}
WITH ranked AS (
    SELECT nombre, salario, departamento_id,
           ROW_NUMBER() OVER (PARTITION BY departamento_id ORDER BY salario DESC) AS rn
    FROM empleados
)
SELECT * FROM ranked WHERE rn <= 3;
        \end{lstlisting}
    \end{block}
\end{frame}

% ============================================================================
% SECCIÓN 3: LIMPIEZA Y TRANSFORMACIÓN DE DATOS
% ============================================================================
\section{Limpieza y Transformación de Datos}

\begin{frame}[fragile]{Manejo de Valores Nulos}
    \begin{block}{\texttt{COALESCE} y \texttt{NULLIF}}
        \begin{lstlisting}
-- Reemplazar nulos
SELECT nombre, COALESCE(email, 'sin-email@empresa.com') AS email
FROM clientes;

-- Evitar división por cero
SELECT total / NULLIF(cantidad, 0) FROM ventas;
        \end{lstlisting}
    \end{block}
\end{frame}

\begin{frame}[fragile]{Funciones de Texto}
    \begin{block}{Manipulación de cadenas}
        \begin{lstlisting}
SELECT UPPER(nombre), LOWER(email), INITCAP(TRIM(nombre))
FROM clientes;

-- Concatenar
SELECT nombre || ' (' || email || ')' FROM clientes;

-- Reemplazar
SELECT REPLACE(telefono, '-', '') FROM clientes;
        \end{lstlisting}
    \end{block}
\end{frame}

\begin{frame}[fragile]{Expresiones Regulares}
    \begin{block}{Extraer dominio de email}
        \begin{lstlisting}
-- PostgreSQL
SELECT email, SUBSTRING(email FROM '@(.*)$') AS dominio
FROM clientes;

-- Con REGEXP_REPLACE
SELECT REGEXP_REPLACE(email, '.*@', '') AS dominio
FROM clientes;
        \end{lstlisting}
    \end{block}
    \begin{block}{Limpiar teléfono (solo dígitos)}
        \begin{lstlisting}
SELECT REGEXP_REPLACE(telefono, '[^0-9]', '', 'g') AS telefono_limpio
FROM clientes;
        \end{lstlisting}
    \end{block}
\end{frame}

\begin{frame}[fragile]{Funciones de Fecha y Hora}
    \begin{block}{Extraer partes, truncar, convertir}
        \begin{lstlisting}
-- Extraer año
SELECT EXTRACT(YEAR FROM fecha) FROM ventas;

-- Truncar a mes
SELECT DATE_TRUNC('month', fecha) FROM ventas;

-- Convertir texto a fecha
SELECT TO_DATE('15/03/2025', 'DD/MM/YYYY') AS fecha;
        \end{lstlisting}
    \end{block}
\end{frame}

\begin{frame}[fragile]{Ejercicio Resuelto: Limpieza de Datos de Clientes}
    \textbf{Enunciado:} Tabla \texttt{clientes\_raw} con columnas: \texttt{id}, \texttt{nombre} (espacios y mayúsculas inconsistentes), \texttt{email} (nulos), \texttt{fecha\_registro} (texto 'DD/MM/YYYY'). Limpiar:
    \begin{enumerate}
        \item Nombre en formato título, sin espacios extras.
        \item Email: si es nulo, asignar 'pendiente@mail.com'.
        \item Fecha convertida a tipo DATE.
    \end{enumerate}
    \begin{block}{Solución}
        \begin{lstlisting}
SELECT id,
       INITCAP(TRIM(nombre)) AS nombre_limpio,
       COALESCE(email, 'pendiente@mail.com') AS email,
       TO_DATE(fecha_registro, 'DD/MM/YYYY') AS fecha_registro_date
FROM clientes_raw;
        \end{lstlisting}
    \end{block}
\end{frame}

% ============================================================================
% SECCIÓN 4: TÉCNICAS AVANZADAS PARA ETL
% ============================================================================
\section{Técnicas Avanzadas para ETL}

\begin{frame}[fragile]{Tablas Temporales}
    \begin{block}{Crear y usar tablas temporales}
        \begin{lstlisting}
CREATE TEMP TABLE ventas_trimestre AS
SELECT *
FROM ventas
WHERE fecha >= DATE_TRUNC('quarter', CURRENT_DATE) - INTERVAL '3 months';

SELECT * FROM ventas_trimestre;
        \end{lstlisting}
    \end{block}
    \begin{itemize}
        \item Existen solo durante la sesión.
        \item Útiles para resultados intermedios.
    \end{itemize}
\end{frame}

\begin{frame}[fragile]{CTEs y Subconsultas Correlacionadas}
    \begin{block}{CTE para mejor legibilidad}
        \begin{lstlisting}
WITH promedio_depto AS (
    SELECT departamento_id, AVG(salario) AS salario_prom
    FROM empleados GROUP BY departamento_id
)
SELECT e.nombre, e.salario, pd.salario_prom
FROM empleados e
JOIN promedio_depto pd ON e.departamento_id = pd.departamento_id
WHERE e.salario > pd.salario_prom;
        \end{lstlisting}
    \end{block}
\end{frame}

\begin{frame}[fragile]{Condicionales con CASE}
    \begin{block}{Categorizar salarios}
        \begin{lstlisting}
SELECT nombre, salario,
       CASE
           WHEN salario < 2000 THEN 'Bajo'
           WHEN salario < 4000 THEN 'Medio'
           ELSE 'Alto'
       END AS categoria_salarial
FROM empleados;
        \end{lstlisting}
    \end{block}
\end{frame}

\begin{frame}[fragile]{Pivotes con CASE y Agregación Condicional}
    \begin{block}{Ventas mensuales por producto}
        \begin{lstlisting}
SELECT producto,
       SUM(CASE WHEN EXTRACT(month FROM fecha) = 1 THEN monto ELSE 0 END) AS enero,
       SUM(CASE WHEN EXTRACT(month FROM fecha) = 2 THEN monto ELSE 0 END) AS febrero,
       SUM(monto) AS total
FROM ventas
GROUP BY producto;
        \end{lstlisting}
    \end{block}
    \begin{block}{Usando FILTER (PostgreSQL)}
        \begin{lstlisting}
SELECT producto,
       SUM(monto) FILTER (WHERE EXTRACT(month FROM fecha) = 1) AS enero,
       SUM(monto) FILTER (WHERE EXTRACT(month FROM fecha) = 2) AS febrero
FROM ventas
GROUP BY producto;
        \end{lstlisting}
    \end{block}
\end{frame}

\begin{frame}[fragile]{Manejo de Duplicados con ROW\_NUMBER}
    \begin{block}{Eliminar duplicados por email}
        \begin{lstlisting}
WITH duplicados AS (
    SELECT *, ROW_NUMBER() OVER (PARTITION BY email ORDER BY id) AS rn
    FROM clientes
)
DELETE FROM clientes
WHERE id IN (SELECT id FROM duplicados WHERE rn > 1);
        \end{lstlisting}
    \end{block}
\end{frame}

\begin{frame}[fragile]{Ejercicio Resuelto: Transformación Compleja}
    \textbf{Enunciado:} Con tablas \texttt{ventas} y \texttt{productos}, se pide:
    \begin{enumerate}
        \item Crear tabla temporal con ventas del último trimestre.
        \item Pivot: total por categoría y mes.
        \item Clasificar ventas como 'baja', 'media', 'alta'.
        \item Mostrar solo categorías con total > 1000.
    \end{enumerate}
    \begin{block}{Solución (parcial)}
        \begin{lstlisting}
-- 1. Tabla temporal
CREATE TEMP TABLE ventas_trimestre AS
SELECT v.*, p.categoria
FROM ventas v
JOIN productos p ON v.producto = p.id
WHERE v.fecha >= DATE_TRUNC('quarter', CURRENT_DATE) - INTERVAL '3 months';

-- 2. Pivot por mes
SELECT categoria,
       SUM(CASE WHEN EXTRACT(month FROM fecha) = 1 THEN monto ELSE 0 END) AS enero,
       SUM(CASE WHEN EXTRACT(month FROM fecha) = 2 THEN monto ELSE 0 END) AS febrero,
       SUM(monto) AS total
FROM ventas_trimestre
GROUP BY categoria
HAVING SUM(monto) > 1000;
        \end{lstlisting}
    \end{block}
\end{frame}

% ============================================================================
% SECCIÓN 5: EJERCICIOS ADICIONALES (DEL SOLUCIONARIO)
% ============================================================================
\section{Ejercicios Adicionales}

\begin{frame}[fragile]{Ejercicio 1: Series Temporales}
    \textbf{Enunciado:} Calcular la media móvil de 7 días de las ventas diarias.
    \begin{block}{Solución}
        \begin{lstlisting}
SELECT fecha, monto,
       AVG(monto) OVER (ORDER BY fecha ROWS BETWEEN 6 PRECEDING AND CURRENT ROW) AS media_movil_7
FROM ventas_diarias;
        \end{lstlisting}
    \end{block}
\end{frame}

\begin{frame}[fragile]{Ejercicio 2: Extraer Código Postal}
    \textbf{Enunciado:} De una dirección como 'Calle Falsa 123, 28080 Madrid', extraer el código postal (5 dígitos).
    \begin{block}{Solución}
        \begin{lstlisting}
SELECT direccion,
       SUBSTRING(direccion FROM '[0-9]{5}') AS codigo_postal
FROM clientes;
        \end{lstlisting}
    \end{block}
\end{frame}

\begin{frame}[fragile]{Ejercicio 3: Agregación Condicional}
    \textbf{Enunciado:} Contar empleados por departamento y cuántos tienen salario > 3000.
    \begin{block}{Solución}
        \begin{lstlisting}
SELECT departamento_id,
       COUNT(*) AS total,
       COUNT(*) FILTER (WHERE salario > 3000) AS altos_salarios
FROM empleados
GROUP BY departamento_id;
        \end{lstlisting}
    \end{block}
\end{frame}

\begin{frame}[fragile]{Ejercicio 4: CTE Recursiva para Jerarquía}
    \textbf{Enunciado:} Con tabla \texttt{empleados(id, nombre, jefe\_id)}, mostrar nivel jerárquico y camino desde el CEO.
    \begin{block}{Solución}
        \begin{lstlisting}
WITH RECURSIVE jerarquia AS (
    SELECT id, nombre, jefe_id, 1 AS nivel, nombre AS path
    FROM empleados WHERE jefe_id IS NULL
    UNION ALL
    SELECT e.id, e.nombre, e.jefe_id, j.nivel + 1,
           j.path || ' > ' || e.nombre
    FROM empleados e
    JOIN jerarquia j ON e.jefe_id = j.id
)
SELECT id, nombre, nivel, path
FROM jerarquia
ORDER BY nivel, nombre;
        \end{lstlisting}
    \end{block}
\end{frame}

% ============================================================================
% SECCIÓN 6: GLOSARIO
% ============================================================================
\section{Glosario}

\begin{frame}{Términos Clave}
    \begin{description}
        \item[Función de agregación] Opera sobre un conjunto y devuelve un único valor.
        \item[GROUP BY] Agrupa filas con valores comunes.
        \item[HAVING] Filtra grupos después de la agregación.
        \item[ROLLUP / CUBE / GROUPING SETS] Extensiones para subtotales.
        \item[Función de ventana] Calcula sobre un marco de filas sin colapsar.
        \item[OVER] Define la ventana (partición, orden, marco).
        \item[ROW\_NUMBER, RANK, DENSE\_RANK] Funciones de ranking.
        \item[LAG / LEAD] Acceden a filas anteriores/posteriores.
        \item[COALESCE] Primer valor no nulo.
        \item[NULLIF] Compara y devuelve NULL si iguales.
        \item[Expresión regular] Patrón para búsqueda en texto.
        \item[CTE] Common Table Expression (WITH).
        \item[Tabla temporal] Tabla efímera para una sesión.
        \item[CASE] Expresión condicional.
        \item[Pivot] Transformar filas en columnas.
    \end{description}
\end{frame}

% --- Diapositiva final ---
\begin{frame}
    \centering

    
    \vspace{1cm}
    \large ¡Gracias!
\end{frame}

\end{document}